\section{Validate the reload and run its display intervals}
The earlier caller checked that reload returned the expected byte count. A player
needs a stronger boundary: validate the actual loaded bytes before using any old
directory or changing the display. Implement \texttt{video\_play} with immutable
file base in \texttt{\$a0}, exact loaded byte count in \texttt{\$a1}, display-image
address in \texttt{\$a2} and a native 96-byte context in \texttt{\$a3}. Source
bytes may be unaligned; display and scratch are four-byte aligned, directory and
context eight-byte aligned. The caller supplies disjoint valid file, display,
directory, scratch, context and private-stack storage. Keep the accepted file and
its new directory immutable throughout this call.

\begin{table}[htbp]
\centering
\caption{Native playback context, with eight-byte fields.}
\label{tab:video-playback-context}
\begin{tabular}{r p{1.5in} p{2.8in}}
\toprule
Offset & Field & Meaning \\
\midrule
0 & Display capacity & At least 16\,384 bytes. \\
8, 16 & Directory address, capacity & Native offsets; enough for the admitted frame count. \\
24, 32 & Scratch address, capacity & Separate working image; at least 16\,384 bytes. \\
40 & Start frame & Unsigned index below the admitted frame count. \\
48 & Accepted frames & Final output, number in the played suffix. \\
56 & Late intervals & Final output, intervals ending after their scheduled deadline. \\
64 & Maximum lateness & Final output, unsigned milliseconds. \\
72 & Logical duration & Final output, suffix interval count times duration. \\
80--95 & Retained suffix & Leave unchanged. \\
\bottomrule
\end{tabular}
\end{table}

Index the whole loaded file before reading the clock or accessing display pixels.
Directory and scratch may change during the call. Seek once to reconstruct the
requested start independently of old displayed pixels, then apply successive
records directly to that accepted image. The bitmap display tool can observe this
image storage when connected to its actual address. No tool connection is required
for the headless pixel checks, and a completed memory store does not establish
when an IDE repaint becomes visible.

Read the initial anchor after validation and before the first seek. Add one duration
to the scheduled deadline before each frame, rejecting unsigned overflow before
presenting that interval. After each accepted frame, read time. Reject an observed
backward change; this player does not re-anchor. If time is below deadline, request
only the positive remainder with service thirty-two, then read time again. Repeat
until deadline is reached. Equal time is on time; a greater value records lateness
once for that interval. Advance the original schedule without dropping frames.
Every frame, including an unchanged or final frame, receives its interval.

This defines the timing protocol for the exercise: one anchor reading, a reading
after each accepted frame, and a further reading after every requested sleep.
There are no other clock requests. It makes the service boundary observable while
allowing different instruction sequences, stack layouts, copies and valid file
representations. Deadlines and observed times use complete unsigned 64-bit values;
all admitted differences fit the stated operation's range. Real clock and sleep
services still provide no guarantee of exact host or UI presentation time.

\begin{table}[htbp]
\centering
\caption{One interval anchored at 1000 with duration forty, in milliseconds.}
\label{tab:video-playback-timelines}
\begin{tabular}{l p{1.5in} p{2.1in}}
\toprule
Observed times after anchor & Requests & Interval result \\
\midrule
1040 & None & Exact deadline, zero lateness. \\
1010, 1040 & Sleep 30 & Remaining time only. \\
1010, 1035, 1040 & Sleep 30, sleep 5 & Early wakeup checked again. \\
1030, 1047 & Sleep 10 & Seven milliseconds late. \\
1010, 1009 & Sleep 30 & Backward observed time: failure. \\
\bottomrule
\end{tabular}
\end{table}

Success returns zero and the accepted frame count. Only complete success publishes
the four final context outputs. Completed failure returns one and the actual count
already accepted; an initial rejection returns count zero. Final context fields
remain unchanged on failure, but earlier accepted display pixels are not rolled
back. A failure after presenting the second frame therefore returns two and leaves
that second image visible. Development status four carries the private progress counter. While the counter
region is still unfinished, an already changed display may coexist with a zero
counter; the complete routine supplies the accepted-count guarantee.
Retain saved registers, frame pointer and original stack pointer, with sixteen-byte
stack alignment across nested calls. The illustrated private frame is 352 bytes,
including separate index, seek and record contexts; another layout is permitted.

Attach your file-codec working source with \texttt{attach.py --source file-work.asm
--output playback-work.asm} from \texttt{packages/034-lossless-video/construction/video-playback-construction}.
The earlier instructions remain intact under \texttt{file\_codec\_example\_main}.
The new caller encodes, checks, seeks, stores and reloads, then invokes the player,
which validates that reload afresh. Prepare observations with \texttt{prepare.py
s102-1 --source playback-work.asm --output observe.asm}. Preparation changes input
data and native controls, retaining your instructions and refusing existing output
files. The first two regions build the callable \texttt{video\_now} helper; guided
checks can call it directly before the player reaches its anchor region.

The grading clock uses \texttt{provideTiming} to supply explicit readings and record
sleep requests without waiting on the host. \texttt{runUntilSyscall} stops before
a selected service so each accepted frame can be inspected without depending on a
private student label. These are deterministic arithmetic and pixel checks. In an
interactive IDE run, inspect the corresponding actual clock values and frame words;
scripted milliseconds are not a claim about the host's speed. Actual annotated IDE
captures remain a separate verification step.
